\subsection{Fixed-target metric realization and its sharp dimensional scope}
\label{sphere-metric-limit-note}

This note concerns the codomains of the three based maps
\[
 q_2^{\wedge 2}:X^{\wedge 2}\cong S^{12}
   \longrightarrow Z_{12}:=\mathcal A^{\wedge 2},
\]
\[
 q_{1,\mathrm{cusp}}^{\wedge 4}:X^{\wedge 4}\cong S^{24}
   \longrightarrow Z_{24}^{\mathcal B}:=\mathcal B^{\wedge 4},
 \qquad
 q_1^{\wedge 4}:X^{\wedge 4}\cong S^{24}
   \longrightarrow Z_{24}^{\mathcal C}:=\mathcal C^{\wedge 4}.
\]
In particular, \(S^{12}\) and \(S^{24}\) are the sources, not the spaces
in which the metric limits occur.  We construct one fixed compatible
length metric on each codomain which realizes the native
four-dimensional flat limits on the target slices.  The realization is
sharp: after the same rescaling, the full ambient limits have dimensions
\(8\) and \(12\), and the collapsed smash base point goes to infinity.

Put
\[
 r_T=e^{-2\pi T},\qquad
 J=\begin{pmatrix}0&1\\-1&0\end{pmatrix},\qquad
 Y_T=TJ+E_T,
\]
where \(E_T\) is the bounded matrix in the radial cusp expansion.  Choose
and fix \(T_0\) and \(C>0\) so that
\(\|E_T\|_{\mathrm{op}}\leq C\) for every \(T\geq T_0\), in the fixed
Euclidean coordinates used below.  This notation avoids confusing the
error matrix with the quotient circle \(C_T\) used below.  On the
fixed torus
\[
 \overline T^2=(\overline\Lambda\otimes\mathbb R)/\overline\Lambda
              \cong\mathbb R^2/\mathbb Z^2
\]
let
\[
 h_T:=\iota_T^*g_{Q_T}=dx^{\mathsf T}Y_T^{\mathsf T}Y_Tdx.
\]
Thus
\begin{equation}
 T^{-2}h_T\longrightarrow h_\infty:=dx^{\mathsf T}dx
 \label{sphere-metric-normalized-Q-limit}
\end{equation}
in every \(C^k\) norm in the fixed flat coordinates.  Indeed, the cusp
coefficients are holomorphic in \(t=e^{-2\pi T}\), so their \(T\)-derivatives
are exponentially small, while \(T^{-1}Y_T=J+T^{-1}E_T\).

Use normalized angular coordinates
\[
 u,\theta\in\mathbb R/\mathbb Z;
\]
\(u\) is the argument of the base coordinate and \(\theta\) is the value of
the invariant character.  Increase this same \(T_0\) once more, if
necessary, so that all radial cusp formulas used below hold for
\(T\geq T_0\) and \(e^{-2\pi T_0}<\epsilon\), and put
\[
 r_0=e^{-2\pi T_0}.
\]
For \(0<r\leq r_0\), write
\[
 T(r):=-\frac{1}{2\pi}\log r.
\]

\paragraph{Target-slice notation.}
We use the embeddings
\(j_T,k_T^{\mathcal B},k_T^{\mathcal C}\) defined in
\cref{cusp-cone-target-slice-embeddings}; the normalized angle \(u=0\) used
below is the same positive radial base angle denoted by \(1\in S^1\) there.

\begin{theorem}[Topology-forced collapse for every fixed compatible metric]
\label{sphere-metric-uniform-collapse}
Let \(D_{12},D_{24}^{\mathcal B},D_{24}^{\mathcal C}\) be arbitrary fixed
metrics compatible with the respective compact target topologies, and let
\(*\) denote the corresponding smash base point.  Then
\begin{align}
 d_H^{D_{12}}\bigl(j_T(Q_T^2),\{*\}\bigr)&\longrightarrow0,\\
 d_H^{D_{24}^{\mathcal B}}
  \bigl(k_T^{\mathcal B}(C_T^4),\{*\}\bigr)&\longrightarrow0,\\
 d_H^{D_{24}^{\mathcal C}}
  \bigl(k_T^{\mathcal C}(C_T^4),\{*\}\bigr)&\longrightarrow0.
 \label{sphere-metric-Hausdorff-collapse}
\end{align}
Consequently, with the metrics restricted from any one fixed ambient
target, all three target slices converge in Gromov--Hausdorff topology to
a point, not to \(\mathbb R^4\).
\end{theorem}

\begin{proof}
We prove the first assertion; the other two are identical.  Let
\[
 G:[0,r_0]^2\times(T^3)^2\longrightarrow Z_{12}
\]
be the product of the two inner-cone quotient maps followed by the smash
quotient.  It is continuous and
\(G(0,0,z_1,z_2)=*\) for every \((z_1,z_2)\in(T^3)^2\).  Since the domain
is compact, uniform continuity with respect to \(D_{12}\) gives
\[
 \sup_{(z_1,z_2)\in(T^3)^2}
 D_{12}\bigl(G(r,r,z_1,z_2),*\bigr)\longrightarrow0.
\]
The set \(j_T(Q_T^2)\) is contained in such a cross-section, which proves
the first Hausdorff convergence.  Replacing \(T^3\) by \(T^2\) and taking
four factors proves the other two.  Hausdorff convergence to a singleton
implies Gromov--Hausdorff convergence to that singleton.
\end{proof}

The preceding theorem is unrescaled.  It does not prevent a blow-up of one
fixed metric.  The following construction gives the strongest natural
positive statement.

\begin{theorem}[One fixed cusp-adapted metric realizes the vertical
\(\mathbb R^4\) limits]
\label{sphere-metric-fixed-realization}
There are fixed compatible length metrics
\[
 D_{12}\text{ on }Z_{12},\qquad
 D_{24}^{\mathcal B}\text{ on }Z_{24}^{\mathcal B},\qquad
 D_{24}^{\mathcal C}\text{ on }Z_{24}^{\mathcal C}
\]
with the following properties.

On the inner cone of \(\mathcal A\), the metric is induced by
\begin{equation}
 g_{\mathcal A}
 =dr^2+r^2du^2+\frac{r^2}{T(r)^2}h_{T(r)}.
 \label{sphere-metric-A-fixed-metric}
\end{equation}
On the cusp cone of either \(\mathcal B\) or \(\mathcal C\), it is induced
by
\begin{equation}
 g_{\mathcal D}=dr^2+r^2(du^2+d\theta^2),
 \qquad\mathcal D\in\{\mathcal B,\mathcal C\}.
 \label{sphere-metric-BC-fixed-metric}
\end{equation}
The metrics on the smash powers are the metric quotients of the product
length metrics after the closed fat wedge is collapsed.  Write
\(g_{Z_{12}}\) and \(g_{Z_{24}^{\mathcal D}}\) for the resulting smooth
Riemannian tensors on the open regular strata.

Let
\[
 p_T^{12}:=j_T(0,0),\qquad
 p_T^{24,\mathcal D}:=k_T^{\mathcal D}(0,0,0,0).
\]
Put
\[
 \Sigma_T^{12}:=j_T(Q_T^2),\qquad
 \Sigma_T^{24,\mathcal D}:=k_T^{\mathcal D}(C_T^4).
\]
Let \(d_{T,12}^{\mathrm{slice}}\) and
\(d_{T,24,\mathcal D}^{\mathrm{slice}}\) be the intrinsic Riemannian path
metrics of the first fundamental forms on these embedded slices.  These
must be distinguished from the restricted ambient distance functions
\[
 d_{T,12}^{\mathrm{amb}}
 :=D_{12}|_{\Sigma_T^{12}\times\Sigma_T^{12}},\qquad
 d_{T,24,\mathcal D}^{\mathrm{amb}}
 :=D_{24}^{\mathcal D}|_{\Sigma_T^{24,\mathcal D}
                         \times\Sigma_T^{24,\mathcal D}}.
\]
In general the latter can be smaller, because an ambient path is allowed
to leave the slice.

Then the induced first fundamental forms on the target slices satisfy the
exact identities
\begin{equation}
 \left(\frac{T}{r_T}\right)^2
  j_T^*g_{Z_{12}}
 =g_{Q_T}\oplus g_{Q_T},
 \label{sphere-metric-A-slice-exact}
\end{equation}
and
\begin{equation}
 \left(\frac{\ell_T}{r_T}\right)^2
  (k_T^{\mathcal D})^*g_{Z_{24}^{\mathcal D}}
 =g_{C_T}^{\oplus4}.
 \label{sphere-metric-BC-slice-exact}
\end{equation}
Consequently,
\begin{align}
 \left(\Sigma_T^{12},\frac{T}{r_T}d_{T,12}^{\mathrm{slice}},
       p_T^{12}\right)
 &\longrightarrow(\mathbb R^4,0),\\
 \left(\Sigma_T^{24,\mathcal D},
       \frac{\ell_T}{r_T}d_{T,24,\mathcal D}^{\mathrm{slice}},
       p_T^{24,\mathcal D}\right)
 &\longrightarrow(\mathbb R^4,0)
 \label{sphere-metric-slice-fixed-blowups}
\end{align}
in pointed local-isometric and pointed Gromov--Hausdorff topology, and
smoothly on the exhausting Euclidean slice balls.

The full ambient pointed blow-ups are instead
\begin{align}
 \left(Z_{12},\frac{T}{r_T}D_{12},p_T^{12}\right)
 &\longrightarrow(\mathbb R^8,0),\\
 \left(Z_{24}^{\mathcal D},
       \frac{\ell_T}{r_T}D_{24}^{\mathcal D},
       p_T^{24,\mathcal D}\right)
 &\longrightarrow(\mathbb R^{12},0).
 \label{sphere-metric-ambient-fixed-blowups}
\end{align}
as pointed metric spaces in pointed Gromov--Hausdorff topology.  More
precisely, for every fixed radius the corresponding rescaled target ball is
eventually contained in an open regular-stratum coordinate neighborhood,
and the pulled-back Riemannian tensors converge there in every \(C^k\) norm.
Thus smooth convergence is asserted on an exhaustion by regular-stratum
neighborhoods, not globally on the singular compact targets.  Under these
limits, the slice \(\mathbb R^4\) is respectively the coordinate subspace
obtained by fixing the two radial and two base-angle coordinates, or by
fixing the four radial and four base-angle coordinates.
\end{theorem}

\begin{proof}
Equation \eqref{sphere-metric-normalized-Q-limit} shows that the metrics on
the compact cross-sections of \eqref{sphere-metric-A-fixed-metric} remain
uniformly bounded and converge to a positive definite flat metric.
Therefore the cross-section diameter is \(O(r)\).  Moreover the radial
coordinate is \(1\)-Lipschitz and a radial segment from level \(r\) to the
vertex has length \(r\).  The length completion consequently adds exactly
one point at \(r=0\), and its topology is the prescribed cone topology.
The same assertions are immediate for
\eqref{sphere-metric-BC-fixed-metric}.  Interpolation on a closed annulus
away from the vertex extends these metrics to fixed compatible length
metrics on the compact spaces \(\mathcal A,\mathcal B,\mathcal C\).

If a compact metric space \((E,d)\) and a closed subset \(W\subset E\) are
given, the metric quotient which collapses \(W\) to one point is
\[
 d_{E/W}([x],[y])
 =\min\{d(x,y),d(x,W)+d(y,W)\}.
\]
Applying this to the fat wedges in the product length spaces gives the
claimed fixed compatible metrics on the smash targets.

On the \(r_T\)-slice of \(\mathcal A\), with \(u=0\), the first fundamental
form is
\[
 \frac{r_T^2}{T^2}h_T.
\]
The definition \(h_T=\iota_T^*g_{Q_T}\) proves
\eqref{sphere-metric-A-slice-exact}.  On the \(r_T\)-slice of
\(\mathcal B\) or \(\mathcal C\), with \(u=0\), the form is
\(r_T^2d\theta^2\).  Under
\(\vartheta_T:[\theta]\mapsto[\ell_T\theta]\), the native quotient-circle
metric is \(\ell_T^2d\theta^2\).  This proves
\eqref{sphere-metric-BC-slice-exact}.  The native systole estimates
\[
 \operatorname{sys}(Q_T^2)\geq T-C,
 \qquad
 \operatorname{sys}(C_T^4)=\ell_T
\]
then prove \eqref{sphere-metric-slice-fixed-blowups}: every prescribed
radius ball is eventually exactly a Euclidean four-ball.

It remains to compute the ambient limits.  In one \(\mathcal A\)-factor,
put
\[
 \lambda_T=\frac{T}{r_T},\qquad
 s=\lambda_T(r-r_T),\qquad
 \eta=Tu,\qquad z=Y_Tx
\]
on the universal cover of the torus directions.  On every fixed bounded
set in \((s,\eta,z)\), one has
\[
 \frac r{r_T}=1+\frac sT,
 \qquad
 T(r)=T+O(T^{-1}),
 \qquad
 Y_{T(r)}Y_T^{-1}=I+O(T^{-2}),
\]
with the same estimates after any fixed number of derivatives.  Thus
\(\lambda_T^2g_{\mathcal A}\) converges in these coordinates to
\[
 ds^2+d\eta^2+dz_1^2+dz_2^2.
\]
The inward radial boundary is at distance \(T\), the \(u\)-period is \(T\),
and the \(z\)-period lattice is \(Y_T\mathbb Z^2\), whose systole is at
least \(T-C\).  The outward edge of the fixed cusp chart is also at scaled
distance
\[
 \frac{T}{r_T}(r_0-r_T)
 =T\left(\frac{r_0}{r_T}-1\right)\longrightarrow\infty.
\]
Thus the inward vertex, all periods, and the outward edge recede to
infinity.  A single factor converges on exhausting regular-stratum
coordinate neighborhoods to \(\mathbb R^4\), and two factors converge to
\(\mathbb R^8\).

For one \(\mathcal B\)- or \(\mathcal C\)-factor, put
\[
 \lambda_T=\frac{\ell_T}{r_T},\qquad
 s=\lambda_T(r-r_T),\qquad
 \eta=\ell_Tu,\qquad \zeta=\ell_T\theta.
\]
Then
\[
 \lambda_T^2g_{\mathcal D}
 =ds^2+\left(\frac r{r_T}\right)^2
       (d\eta^2+d\zeta^2),
\]
which converges smoothly on bounded coordinate sets to the Euclidean metric
on \(\mathbb R^3\).  The two circle periods and the inward radial boundary
are all at scaled distance \(\ell_T\to\infty\), while the outward edge is
at scaled distance
\[
 \frac{\ell_T}{r_T}(r_0-r_T)
 =\ell_T\left(\frac{r_0}{r_T}-1\right)\longrightarrow\infty.
\]
Four factors therefore converge on exhausting regular-stratum coordinate
neighborhoods to \(\mathbb R^{12}\).

Finally, a level-\(r_T\) factor point has distance exactly \(r_T\) from its
cone vertex.  Indeed, the radial segment has that length.  A competing
curve contained in the inner chart has length at least \(r_T\), while one
which leaves has inner-chart portions of total length at least
\((r_0-r_T)+r_0>r_T\).  The Riemannian-\(\ell^2\)-product distance from a
synchronized point to the fat wedge is the minimum of the factor
distances, and hence is also \(r_T\).

For the chosen scale put
\[
 b_T=\lambda_Tr_T,\qquad a_T=\lambda_T(r_0-r_T),
\]
and let \(\sigma_T\) be the least scaled period in a factor.  Thus
\(\sigma_T\geq T-C\) in the \(\mathcal A\) calculation and
\(\sigma_T=\ell_T\) in the \(\mathcal B,\mathcal C\) calculation.  The
preceding computations give
\[
 a_T\longrightarrow\infty,\qquad
 b_T\longrightarrow\infty,\qquad
 \sigma_T\longrightarrow\infty.
\]
Fix \(R\).  A curve of rescaled length at most \(R\) from \(p_T\) meets
neither the fat wedge nor \(r=r_0\) for all large \(T\), so the
radius-\(R\) target ball lies in the open regular stratum.  If \(x,y\) lie
in that ball, a competitor through the fat wedge has length at least
\[
 (b_T-R)+(b_T-R)>2R,
\]
and a competitor which exits the inner chart and returns has length at
least
\[
 2(a_T-R)>2R.
\]
Every remaining competitor of length at most \(2R\) stays in
\(r/r_T=1+o_R(1)\).  After lifting it to the angular and torus covers, a
nonzero deck translation has length at least
\((1-o_R(1))(\sigma_T-2R)>2R\) for all large \(T\).  Since the path through
\(p_T\) has length at most \(2R\), none of these three kinds of global
shortcut computes the distance between \(x\) and \(y\).  Thus the metric
on every fixed target ball is computed entirely in one lifted
regular-stratum chart.  The \(C^\infty\) convergence of the pulled-back
tensors on these exhausting charts implies the pointed
Gromov--Hausdorff limits in
\eqref{sphere-metric-ambient-fixed-blowups}.
\end{proof}

\begin{theorem}[Complete intrinsic scalar-rescaling trichotomy]
\label{sphere-metric-scaling-trichotomy}
For the fixed metrics of
\cref{sphere-metric-fixed-realization}, let \(\lambda_T>0\) be arbitrary
and put
\[
 \beta_T:=\lambda_Tr_T.
\]
Choose one of the following triples, retaining the indicated ambient
metric:
\[
 \begin{split}
  (\Sigma_T,d_T^{\mathrm{slice}},d_T^{\mathrm{amb}},p_T;D)
  &=(\Sigma_T^{12},d_{T,12}^{\mathrm{slice}},
      d_{T,12}^{\mathrm{amb}},p_T^{12};D_{12}),\\
  (\Sigma_T,d_T^{\mathrm{slice}},d_T^{\mathrm{amb}},p_T;D)
  &=(\Sigma_T^{24,\mathcal D},
      d_{T,24,\mathcal D}^{\mathrm{slice}},
      d_{T,24,\mathcal D}^{\mathrm{amb}},
      p_T^{24,\mathcal D};D_{24}^{\mathcal D}),
      \qquad \mathcal D\in\{\mathcal B,\mathcal C\}.
 \end{split}
\]
After passing to a subsequence so that \(\beta_T\) has a limit in
\([0,\infty]\), the intrinsic pointed slice limit is
\begin{equation}
 (\Sigma_T,\lambda_Td_T^{\mathrm{slice}},p_T)
 \longrightarrow
 \begin{cases}
  \{\mathrm{pt}\},&\beta_T\longrightarrow0,\\
  ((\mathbb R/L\mathbb Z)^4,0),
      &\beta_T\longrightarrow L\in(0,\infty),\\
  (\mathbb R^4,0),&\beta_T\longrightarrow\infty.
 \end{cases}
 \label{sphere-metric-scaling-trichotomy-result}
\end{equation}
The convergence in
\eqref{sphere-metric-scaling-trichotomy-result} is pointed
Gromov--Hausdorff in all three cases.  In the finite positive case the flat
metrics converge smoothly under the fixed marked-torus identifications; in
the unrolling case the convergence is pointed local-isometric and smooth on
the exhausting Euclidean balls.
In the middle line the lattice is orthogonally equivalent to the square
lattice \(L\mathbb Z^4\).  At the same time,
\begin{equation}
 \lambda_TD(p_T,*)=\beta_T
 \label{sphere-metric-vertex-distance}
\end{equation}
for every one of the displayed choices.

Only in the unrolling regime \(\beta_T\to\infty\) do we compare the
intrinsic metric with the restricted ambient distance.  In that regime, for
every \(R\in(0,\infty)\),
\begin{equation}
 \sup_{x,y\in
  \overline B_{\lambda_Td_T^{\mathrm{slice}}}(p_T,R)}
 \left|
  \lambda_Td_T^{\mathrm{amb}}(x,y)
  -\lambda_Td_T^{\mathrm{slice}}(x,y)
 \right|\longrightarrow0.
 \label{sphere-metric-intrinsic-ambient-agreement}
\end{equation}
Moreover, there are \(\varepsilon_{T,R}\to0\) such that
\begin{equation}
 \overline B_{\lambda_Td_T^{\mathrm{slice}}}(p_T,R)
 \subset
 \overline B_{\lambda_Td_T^{\mathrm{amb}}}(p_T,R)
 \subset
 \overline B_{\lambda_Td_T^{\mathrm{slice}}}
             (p_T,(1+\varepsilon_{T,R})R).
 \label{sphere-metric-intrinsic-ambient-ball-comparison}
\end{equation}
Thus the restricted ambient metric has the same pointed
\(\mathbb R^4\) limit on these fixed slice balls.  No agreement is asserted
in the regimes \(\beta_T\to0\) or \(\beta_T\to L\in(0,\infty)\), where
radial or fat-wedge shortcuts need not disappear.  Every rescaling which
unrolls the intrinsic slice sends the cone/fat-wedge base point to infinity,
so no such limit is a tangent cone based at that point.

In the same regime \(\beta_T\to\infty\), the corresponding full ambient
bubbles are
\begin{align}
 (Z_{12},\lambda_TD_{12},p_T^{12})
 &\longrightarrow(\mathbb R^8,0),\\
 (Z_{24}^{\mathcal D},\lambda_TD_{24}^{\mathcal D},
  p_T^{24,\mathcal D})
 &\longrightarrow(\mathbb R^{12},0),
 \qquad\mathcal D\in\{\mathcal B,\mathcal C\}.
 \label{sphere-metric-arbitrary-unrolling-ambient-bubbles}
\end{align}
These are pointed Gromov--Hausdorff limits of the singular metric targets.
The convergence of tensors is \(C^\infty\) only under coordinate
identifications on exhausting open regular-stratum neighborhoods; no
global smooth convergence of the singular targets is claimed.
\end{theorem}

\begin{proof}
For the \(Q_T^2\) route, the rescaled period matrix in each two-dimensional
factor is
\[
 \frac{\lambda_Tr_T}{T}Y_T
 =\beta_T\left(J+\frac{E_T}{T}\right).
\]
If \(\beta_T\to L\in(0,\infty)\), this lattice converges to
\(LJ\mathbb Z^2\); two factors give the square four-torus in
\eqref{sphere-metric-scaling-trichotomy-result}.  If
\(\beta_T\to0\), the lengths of a lattice basis tend to zero, so the
diameter tends to zero.  If \(\beta_T\to\infty\), the reverse triangle
inequality gives systole
\[
 \frac{\beta_T}{T}(T-C)
 =\beta_T\left(1-\frac CT\right)\longrightarrow\infty,
\]
so every fixed-radius ball is eventually Euclidean.

For a \(C_T^4\) slice, the four periods after rescaling have the common
length
\[
 \lambda_Tr_T=\beta_T,
\]
which gives the same three intrinsic cases directly.

Consider first one factor.  A curve from a level-\(r_T\) point to the cone
vertex which remains in the inner chart has length at least \(r_T\), because
the \(dr^2\)-term makes \(r\) \(1\)-Lipschitz there.  If it leaves the
chart, it has a first crossing of \(r=r_0\) and, before ending at the
vertex, a last re-entry through \(r=r_0\); its two inner-chart portions
have total length at least
\((r_0-r_T)+r_0>r_T\).  The radial segment has length \(r_T\), so the
factor distance to its based vertex is exactly \(r_T\).  The
Riemannian-\(\ell^2\)-product distance from every synchronized slice point
to the fat wedge is the minimum of these factor distances and hence is
also \(r_T\).  The metric-quotient formula in the preceding proof therefore
gives \eqref{sphere-metric-vertex-distance}.

It remains to prove the restricted-ambient comparison, so assume
\(\beta_T\to\infty\) and fix \(R\).  Two points in
\(\overline B_{\lambda_Td_T^{\mathrm{slice}}}(p_T,R)\) have intrinsic
distance at most \(2R/\lambda_T\).  The slice path is an ambient competitor,
so
\[
 d_T^{\mathrm{amb}}(x,y)\leq d_T^{\mathrm{slice}}(x,y).
\]
The quotient shortcut through the fat wedge has rescaled length
\(2\lambda_Tr_T=2\beta_T\), and hence is longer than \(2R\) for all large
\(T\).  A path which leaves the fixed cusp chart must reach \(r=r_0\).
Its outward cost is controlled by
\begin{equation}
 \lambda_T(r_0-r_T)
 =\beta_T\left(\frac{r_0}{r_T}-1\right)\longrightarrow\infty,
 \label{sphere-metric-outward-shortcut-exclusion}
\end{equation}
so no path of rescaled length at most \(2R\) can use the interpolating
annulus or the global part of a factor.

Consider now an ambient competitor of length at most
\(2R/\lambda_T+o(\lambda_T^{-1})\) which stays in the cusp chart.  In every
factor its radial coordinate satisfies
\[
 |r-r_T|\leq \frac{2R+o(1)}{\lambda_T},
 \qquad
 \frac r{r_T}=1+O_R(\beta_T^{-1}).
\]
On this tube, the horizontal tensors in the directions tangent to the slice
are uniformly \((1+\varepsilon_{T,R})\)-bilipschitz to their values at
\(r=r_T\), where \(\varepsilon_{T,R}\to0\).  For
\(\mathcal B\) and \(\mathcal C\) this follows directly from
\((r/r_T)^2=1+O_R(\beta_T^{-1})\).  For \(\mathcal A\), the same conclusion
follows, without changing the fixed coordinates or the matrix \(Y_T\), from
\[
 \begin{aligned}
  T(r)-T&=O_R(\beta_T^{-1}),\\
  \frac{T(r)}T&=1+O_R((\beta_TT)^{-1}),\\
  Y_{T(r)}Y_T^{-1}&=I+O_R((\beta_TT)^{-1}).
 \end{aligned}
\]
Projecting the competitor to the torus or character directions, and
discarding its radial and base-angle components, produces a path between
the same slice endpoints.  The bilipschitz estimate therefore gives
\[
 d_T^{\mathrm{slice}}(x,y)
 \leq(1+\varepsilon_{T,R})d_T^{\mathrm{amb}}(x,y).
\]
Together with the opposite inequality and
\(\lambda_Td_T^{\mathrm{slice}}(x,y)\leq2R\), this proves
\eqref{sphere-metric-intrinsic-ambient-agreement}.  The first inclusion in
\eqref{sphere-metric-intrinsic-ambient-ball-comparison} follows from
\(d_T^{\mathrm{amb}}\leq d_T^{\mathrm{slice}}\).  For a point \(x\) in
the middle ball, apply the same fat-wedge and outward exclusions to a
near-minimizing ambient path from \(p_T\) to \(x\), whose rescaled length is
at most \(R+o(1)\), and project that path to the slice directions.  This
gives
\[
 \lambda_Td_T^{\mathrm{slice}}(p_T,x)
 \leq(1+\varepsilon_{T,R})R
\]
after enlarging \(\varepsilon_{T,R}\) by a quantity tending to zero, and
proves the second inclusion.

For completeness, the full ambient bubbles at an arbitrary unrolling scale
are obtained without changing any coordinate or orientation.  In one
\(\mathcal A\)-factor use
\[
 s=\lambda_T(r-r_T),\qquad
 \eta=\beta_Tu,\qquad
 z=\frac{\beta_T}{T}Y_Tx.
\]
Writing \(A_T(r)=Y_{T(r)}Y_T^{-1}\), one obtains the exact expression
\[
 \lambda_T^2g_{\mathcal A}
 =ds^2+\left(\frac r{r_T}\right)^2d\eta^2
  +\left(\frac r{r_T}\right)^2
   \left(\frac{T}{T(r)}\right)^2
   dz^{\mathsf T}A_T(r)^{\mathsf T}A_T(r)dz.
\]
On every fixed bounded coordinate set,
\[
 \frac r{r_T}=1+O_R(\beta_T^{-1}),\qquad
 T(r)-T=O_R(\beta_T^{-1}),\qquad
 A_T(r)=I+O_R((\beta_TT)^{-1}),
\]
with the same estimates after every fixed number of derivatives.  Hence
the tensor converges smoothly to the Euclidean metric on \(\mathbb R^4\).
The \(\eta\)-period is \(\beta_T\), and the \(z\)-period lattice is
\[
 \frac{\beta_T}{T}Y_T\mathbb Z^2,
 \qquad
 \operatorname{sys}\left(
   \frac{\beta_T}{T}Y_T\mathbb Z^2\right)
 \geq\beta_T\left(1-\frac CT\right)\longrightarrow\infty.
\]
In one \(\mathcal B\)- or \(\mathcal C\)-factor use
\[
 s=\lambda_T(r-r_T),\qquad
 \eta=\beta_Tu,\qquad \zeta=\beta_T\theta.
\]
Then
\[
 \lambda_T^2g_{\mathcal D}
 =ds^2+\left(\frac r{r_T}\right)^2
       (d\eta^2+d\zeta^2),
\]
and both angular periods are \(\beta_T\to\infty\).

In all three cases the inward boundary and the fat wedge are at scaled
distance \(\beta_T\), while a lift starting at radius \(r_T\) incurs the
one-way radial cost in
\eqref{sphere-metric-outward-shortcut-exclusion} before its first outward
crossing of \(r=r_0\).  Fix \(R\).  For large \(T\), all these quantities
and all nonzero period lengths exceed \(4R\).
If \(x,y\) lie in the ambient rescaled radius-\(R\) ball, their distances
from the fat wedge are at least \(\beta_T-R\); hence a quotient path through
the smash base point has length at least \(2(\beta_T-R)>2R\).  Every
remaining competitor lifts to the complement of the fat wedge in the
product.  If such a lift exits the inner product chart, then some factor
has a first crossing of \(r=r_0\); the radial variation before that crossing
is at least \(\lambda_T(r_0-r_T)-R\), and returning to the lift of \(y\)
incurs the same lower bound.  Its length is therefore greater than \(2R\).
Finally, a path of length at most \(2R\) remains in the tube
\(r/r_T=1+O_R(\beta_T^{-1})\); there every nonzero deck translation has
length tending to infinity, so such a path cannot wrap around an angular
or torus period.  The displayed coordinate neighborhoods therefore exhaust
the pointed balls without quotient or global shortcuts.
Their smooth tensor convergence proves the pointed Gromov--Hausdorff
limits in
\eqref{sphere-metric-arbitrary-unrolling-ambient-bubbles}, with smooth
convergence confined to this regular-stratum exhaustion.
\end{proof}

\begin{corollary}[Unit-volume normalization does not give
\(\mathbb R^4\)]
\label{sphere-metric-unit-volume}
The four-volume of \(Q_T^2\) is \((\det Y_T)^2\sim T^4\).  Its unit-volume
length factor is therefore \((\det Y_T)^{-1/2}\sim T^{-1}\), and its
unit-volume limit is a square flat four-torus, not \(\mathbb R^4\).
Likewise, the four-volume of \(C_T^4\) is \(\ell_T^4\), and its unit-volume
normalization is exactly the square unit four-torus.
\end{corollary}

\begin{proposition}[What vertex tangents and compact asymptotic cones can
mean]
\label{sphere-metric-tangent-scope}
The compact target topology alone determines no metric tangent cone.
For an exact conic length metric built from fixed piecewise-Riemannian
metrics on the links, however, the vertex links have topological and
Hausdorff dimensions \(7\) for \(Z_{12}\) and \(11\) for both
\(Z_{24}\) targets.  The resulting vertex cones have topological and
Hausdorff dimensions \(8\) and \(12\), and are not \(\mathbb R^4\).

More precisely, writing \(L_3=T^3\) and \(L_2=T^2\), one has the based
topological models
\[
 Z_{12}\cong\operatorname{Susp}(L_3*L_3),
 \qquad
 Z_{24}^{\mathcal B}
 \cong\operatorname{Susp}(L_2*L_2*L_2*L_2).
\]
At the suspension base point their conic links are respectively
\(L_3*L_3\) and \(L_2^{*4}\).  Also
\[
 (\mathcal C,c_0)
 \cong
 (\mathcal M/\partial\mathcal M,[\partial\mathcal M]),
 \qquad
 (Z_{24}^{\mathcal C},*)
 \cong
 (\mathcal M^4/\partial(\mathcal M^4),
  [\partial(\mathcal M^4)]),
\]
so an exact conic metric at the collapsed boundary has link
\(\partial(\mathcal M^4)\).  At the collapsed point of
\(Z_{24}^{\mathcal C}\), the complete integral local-homology calculation is
\begin{equation}
 H_i\bigl(Z_{24}^{\mathcal C},
          Z_{24}^{\mathcal C}\setminus\{*\};\mathbb Z\bigr)
 \cong
 \begin{cases}
  \mathbb Z^{\binom4{i-1}},&2\leq i\leq5,\\
  0,&i=6,7,\\
  \mathbb Z^{\binom4{i-8}},&8\leq i\leq12,\\
  0,&\text{otherwise}.
 \end{cases}
 \label{sphere-metric-C-smash-local-homology}
\end{equation}
In particular, its local homology is not concentrated in top degree.  For
the first two targets the top local homology occurs in degree \(8\) and
\(12\), respectively; for the third target
\eqref{sphere-metric-C-smash-local-homology} includes a top
\(\mathbb Z\) in degree \(12\).  None has the local homology of
\(\mathbb R^4\).

Finally, the standard asymptotic cone obtained by blowing down any fixed
compact metric target is a point.  Indeed, if \(R_n\to\infty\), then
\[
 \operatorname{diam}(Z,R_n^{-1}D)
 =R_n^{-1}\operatorname{diam}(Z,D)\longrightarrow0.
\]
Thus the limits in \cref{sphere-metric-fixed-realization} are
moving-base-point blow-ups, often called bubbles.  They are neither
vertex-based tangent cones nor asymptotic cones of the compact targets.
\end{proposition}

\begin{proof}
For compact spaces \(V_i\) with closed subspaces \(A_i\), there is a
canonical homeomorphism
\[
 (V_1/A_1)\wedge\cdots\wedge(V_m/A_m)
 \cong
 \frac{V_1\times\cdots\times V_m}
      {\bigcup_iV_1\times\cdots\times A_i\times\cdots\times V_m}.
\]
Take \(V_i=\operatorname{Cone}(L_i)\) and \(A_i=L_i\), its outer boundary.
The product of \(m\) cones is homeomorphic to the cone on the iterated join
\(L_1*\cdots*L_m\), and its displayed union of faces is the outer link.
Collapsing that link gives its suspension.  This proves the first two
models and the link dimensions
\[
 \dim(T^3*T^3)=3+3+1=7,
 \qquad
 \dim((T^2)^{*4})=4\cdot2+3=11.
\]

We next prove the based homeomorphism for \(\mathcal C\), rather than
identifying the two presentations without a map.  Put
\(L=\partial\mathcal M\), write the attached cone as
\[
 \operatorname{Cone}(L)
 =L\times[0,1]/(L\times\{0\}),
\]
with base \(L\times\{1\}\), and choose a collar
\(\kappa:L\times[0,1]\hookrightarrow\mathcal M\) satisfying
\(\kappa(x,0)=x\).  There is a homeomorphism
\[
 \Phi:\mathcal C\setminus\{c_0\}
       \longrightarrow\operatorname{int}\mathcal M
\]
given on the punctured cone, the inner collar, and the remaining core by
\[
 \Phi([x,t])=\kappa(x,t/e)\quad(0<t\leq1),\qquad
 \Phi(\kappa(x,s))=\kappa(x,e^{s-1})\quad(0\leq s\leq1),
\]
\[
 \Phi(q)=q
 \quad\text{for}\quad
 q\in\mathcal M\setminus\kappa(L\times[0,1)).
\]
The first two formulas agree at the glued points
\([x,1]=\kappa(x,0)\), and the second and third agree at
\(\kappa(x,1)\).  Their collar-coordinate images are respectively
\(L\times(0,e^{-1}]\) and \(L\times[e^{-1},1]\); hence the displayed
piecewise map and its piecewise inverse are continuous.

The compact Hausdorff spaces \(\mathcal C\) and
\(\mathcal M/L\) are both one-point compactifications of
\(\operatorname{int}\mathcal M\): for \(\mathcal C\) use \(\Phi\), and for
\(\mathcal M/L\) use the quotient map on
\(\mathcal M\setminus L=\operatorname{int}\mathcal M\).  Uniqueness of the
one-point compactification therefore extends \(\Phi\) to a based
homeomorphism
\begin{equation}
 (\mathcal C,c_0)
 \xrightarrow{\ \cong\ }
 (\mathcal M/\partial\mathcal M,[\partial\mathcal M]).
 \label{sphere-metric-C-boundary-quotient-homeomorphism}
\end{equation}
Only after establishing
\eqref{sphere-metric-C-boundary-quotient-homeomorphism} do we apply the
quotient identity at the start of the proof.  It gives
\[
 \mathcal C^{\wedge4}
 \cong
 \frac{\mathcal M^4}
 {\displaystyle\bigcup_{a=1}^4
  \mathcal M^{a-1}\times\partial\mathcal M
  \times\mathcal M^{4-a}}
 =\mathcal M^4/\partial(\mathcal M^4),
\]
and carries the smash base point to the collapsed boundary.

Put \(N=\mathcal M^4\).  Rounding its product corners preserves the
homeomorphism type of the pair \((N,\partial N)\) and makes
\(\partial N\) an eleven-dimensional closed manifold.  A boundary collar
descends in \(N/\partial N\) to a cone neighborhood of the collapsed
point, so the local link is exactly \(\partial N\) and
\begin{equation}
 H_i(N/\partial N,(N/\partial N)\setminus\{*\};\mathbb Z)
 \cong\widetilde H_{i-1}(\partial N;\mathbb Z).
 \label{sphere-metric-boundary-link-local-shift}
\end{equation}

The integral homology calculation in the proof of
\cref{cusp-cone-S24-route} for the specific compact orientable
three-manifold \(\mathcal M\) gives
\[
 H_j(\mathcal M;\mathbb Z)\cong
 \begin{cases}
  \mathbb Z,&j=0,1,\\
  0,&\text{otherwise}.
 \end{cases}
\]
All groups are free, so the integral Kunneth theorem gives
\[
 H_j(N;\mathbb Z)\cong
 \begin{cases}
  \mathbb Z^{\binom4j},&0\leq j\leq4,\\
  0,&\text{otherwise}.
 \end{cases}
\]
Poincare--Lefschetz duality, with the product orientation retained, gives
\[
 H_j(N,\partial N;\mathbb Z)
 \cong H^{12-j}(N;\mathbb Z)
 \cong
 \begin{cases}
  \mathbb Z^{\binom4{j-8}},&8\leq j\leq12,\\
  0,&\text{otherwise}.
 \end{cases}
\]
The long exact sequence of \((N,\partial N)\) now has disjoint absolute
and relative degree ranges.  It gives
\[
 \widetilde H_j(\partial N;\mathbb Z)\cong
 \begin{cases}
  \mathbb Z^{\binom4j},&1\leq j\leq4,\\
  0,&j=5,6,\\
  \mathbb Z^{\binom4{j-7}},&7\leq j\leq11,\\
  0,&\text{otherwise}.
 \end{cases}
\]
Substitution in
\eqref{sphere-metric-boundary-link-local-shift} proves
\eqref{sphere-metric-C-smash-local-homology}.

An exact metric cone is invariant under every radial blow-up at its vertex,
and its Hausdorff dimension is one plus that of its fixed
piecewise-Riemannian link.  For the first target,
\[
 H_8\bigl(\operatorname{Cone}(T^3*T^3),T^3*T^3;\mathbb Z\bigr)
 \cong\mathbb Z;
\]
the analogous top local group for the \(\mathcal B\) target is
\(\mathbb Z\) in degree \(12\), while the full \(\mathcal C\) calculation
is \eqref{sphere-metric-C-smash-local-homology}.  Euclidean
\(\mathbb R^4\) has local homology \(\mathbb Z\) only in degree \(4\), so
none of these conic vertex models is \(\mathbb R^4\).  The diameter
calculation proves the final assertion.
\end{proof}

\paragraph{Consequences for the three sphere routes.}
The exact positive statement is that one fixed metric on each compact
codomain contains the native \(\mathbb R^4\) decompactification as a
vertical four-plane in a moving-point Euclidean bubble.  The whole bubbles
are \(\mathbb R^8\) for the \(S^{12}\) route and \(\mathbb R^{12}\) for the
\(S^{24}\) routes.  The factor \(\beta_T=\lambda_Tr_T\) proves that this
dimension discrepancy cannot be removed by a different scalar rescaling
of the same fixed cusp-adapted metric: the slice unrolls only when the
smash base point has already gone to infinite scaled distance.

No metric on a source sphere is supplied by these quotient constructions.
If an independent smooth Riemannian metric is placed on \(S^{12}\) or
\(S^{24}\), its ordinary tangent spaces are \(\mathbb R^{12}\) or
\(\mathbb R^{24}\).  Pulling a codomain metric back through
\(q_2^{\wedge2}\) or \(q_1^{\wedge4}\), on the regular strata where the
maps are differentiable, gives a degenerate quadratic form along the
positive-dimensional quotient kernels, not a Riemannian source metric.

The moving four-sphere carriers
\[
 \mathbb S_{12,T}^4=S_T^2\wedge S_T^2\subset S^6,\qquad
 \mathbb S_{24,T}^4=C_T^{\wedge4}\subset S^8
\]
of \cref{cusp-cone-angle-forgotten-sphere-carriers} carry separately
assigned metrics.  Namely, \(S_T^2\) first receives the quotient length
metric obtained from
\((\overline T^2,h_T)\cong(Q_T,g_{Q_T})\) by collapsing
\(K^{(1)}\subset\overline T^2\), equivalently by collapsing
\(\iota_T(K^{(1)})\subset Q_T\).  Then
\(\mathbb S_{12,T}^4\) receives the quotient length metric of the
Riemannian \(\ell^2\)-product \((S_T^2)^2\) after its fat wedge is
collapsed.  Likewise, \(\mathbb S_{24,T}^4\) receives the quotient length
metric of the Riemannian \(\ell^2\)-product \(C_T^4\) after its fat wedge is
collapsed.  These metrics are assigned directly to the moving carriers;
they are unrescaled intrinsic quotient metrics, not restrictions of
\(D_{12}\) or \(D_{24}^{\mathcal B}\), whose domains are the conical
targets, and no fixed ambient metric on \(S^6\) or \(S^8\) is used to
define them.  Their Euclidean balls are
based at the midpoint classes
\[
 \begin{aligned}
  y_*&=[(1/2,1/2)]\in\overline T^2,\qquad z_T=c(y_*),\\
  \zeta_{12,T}&=[z_T,z_T],\\
  \zeta_{24,T}
  &=[(\ell_T/2,\ell_T/2,\ell_T/2,\ell_T/2)].
 \end{aligned}
\]
not at the images of \(p_T^{12}\) or \(p_T^{24,\mathcal B}\), which belong
to the subsequently collapsed loci.  In particular, the trichotomy in
\cref{sphere-metric-scaling-trichotomy} is the trichotomy of the
uncollapsed intrinsic flat-torus slices; it is not a finite-\(\beta_T\)
description of the singular quotient-metric moving four-spheres.  Indeed,
the restriction of any one compatible metric on the compact sphere
\(S^6\) or \(S^8\) has diameter bounded by the diameter of that sphere.  By
contrast, in the separately assigned moving metrics the midpoint distances
to the collapsed loci are at least \((T-C)/2\) and exactly
\(\ell_T/2\), respectively, and hence tend to infinity.  Thus the assigned
unrescaled metrics cannot themselves be restrictions of one fixed ambient
metric.  This does not preclude a separately proved realization after a
scalar blow-up of another fixed quotient metric.

Finally, the rank obstruction for a full-base torus quotient must retain
its actual metric hypothesis: \emph{quotient systole tends to infinity}
means the quotient norm induced by the unrescaled Euclidean fibre metric.
Without that qualification the assertion is false, since an arbitrary
\(T\)-dependent scalar multiplication can make the systole of any fixed
nonzero quotient tend to infinity.
